\section{Билет 15. Затухающие колебания. Уравнение затухающих колебаний и его решение. Декремент затухания. Логарифмический декремент затухания. Добротность. Апериодические колебания. Демпфер.}

\begin{definition}
Уравнение затухающих колебаний имеет вид:
\begin{equation}
    \ddot{x} + 2\beta \dot{x} + \omega_0^2 x = 0, \label{eq:damped_15}
\end{equation}
где $\beta = r/(2m)$ --- коэффициент затухания, $\omega_0 = \sqrt{k/m}$ --- собственная частота. При малом затухании ($\beta < \omega_0$) решением является:
\begin{equation}
    x(t) = A_0 e^{-\beta t} \cos(\omega t + \alpha), \quad \omega = \sqrt{\omega_0^2 - \beta^2}.
\end{equation}
Амплитуда убывает экспоненциально: $A(t) = A_0 e^{-\beta t}$.
\end{definition}

\begin{definition}
\textbf{Декрементом затухания} $\theta$ называется отношение амплитуд двух последовательных колебаний, отстоящих друг от друга на один период $T$:
\begin{equation}
    \theta = \frac{A(t)}{A(t+T)} = \frac{A_0 e^{-\beta t}}{A_0 e^{-\beta (t+T)}} = e^{\beta T}.
\end{equation}
\textbf{Логарифмическим декрементом затухания} $\lambda$ называется натуральный логарифм декремента:
\begin{equation}
    \lambda = \ln \theta = \beta T. \label{eq:log_decrement}
\end{equation}
Величина $\lambda$ характеризует скорость убывания амплитуды. Через неё коэффициент затухания выражается как $\beta = \lambda / T$.
\end{definition}

\begin{property}[Добротность колебательной системы]
\textbf{Добротностью} $Q$ называют безразмерную величину, характеризующую скорость потери энергии колебательной системой:
\begin{equation}
    Q = \frac{\pi}{\lambda} = \pi N_e, \label{eq:quality_factor}
\end{equation}
где $N_e = \tau/T = 1/(\beta T)$ --- число полных колебаний, совершаемых системой за время релаксации $\tau$. Чем больше добротность, тем медленнее затухают колебания и тем меньше потери энергии за период.
\end{property}

\begin{remark}
\textbf{Апериодические колебания.} При увеличении коэффициента затухания $\beta$ частота $\omega$ уменьшается, а период $T$ растёт. Когда $\beta \to \omega_0$, период обращается в бесконечность. При $\beta \ge \omega_0$ движение становится \textbf{апериодическим}: выведенная из равновесия система медленно возвращается в положение равновесия, не совершая колебаний. Вся начальная энергия расходуется на преодоление сил трения.
\par\noindent\textbf{Демпфер} --- устройство для гашения (демпфирования) или предотвращения нежелательных колебаний в машинах, конструкциях и приборах. Работает на принципе перевода механической энергии колебаний во внутреннюю энергию вязкой среды или тепла (примеры: дверные доводчики, амортизаторы в автомобилях, гасители вибраций в строительстве).
\end{remark}
